\section{Discussion and limitations}\label{sec:discussion}

\paragraph{Reference models.} The energy, wind and perception models are reference models: a P600-class hexacopter
power model, an analytic urban wind and visibility field driven by twelve weather presets, and the Johnson criteria with
Koschmieder contrast attenuation. We report relative differences between policies on the same ground truth because
these differences depend on the coupling, not on the calibration of one vehicle. Absolute endurance, wind loads near
buildings and detection probabilities of a specific drone and camera will differ, and calibrating the models against
flight logs, measured urban wind and camera imagery is the next step. Because every model sits behind one interface,
a calibrated model replaces the reference one for all modules at once, which is the point of sharing it.

\paragraph{Geometry.} Safety is certified against the city's surface model at 2\,m resolution, dilated by 4\,m and
raised by a 10\,m margin. Wires, cranes, thin structures, overhangs and changes after the point cloud was captured are
not in the model, and the margin is not a substitute for onboard sensing. The runtime treats the surface model as one
input to the decision and the separation and contact stages as another; an onboard obstacle-avoidance layer remains
necessary in real operations.

\paragraph{Weather in time.} The environment field represents a forecast as a keyframe transition, and the space-time
precheck uses it as if it were exact. Real forecasts are uncertain in timing and strength; a precheck that plans
against an ensemble of transitions, and a return trigger that re-plans when the forecast changes in flight, are natural
extensions. The precheck also samples the route wind at a few points per leg, which left 4--11\% of sorties meeting
wind beyond the rating on the path they actually flew.

\paragraph{Scope of the evaluation.} Scan plans are scored on their planned trajectories with the perception model
rather than flown in closed loop, and the closed-loop campaigns use eight drones per task; the 1,000-drone measurements
show that the runtime keeps the coupled return logic inside the real-time budget, not that decisions remain equally
good when hundreds of drones interact. Browser access was measured on a laptop, a software renderer and a data-centre
GPU, and the server's cost per additional viewer was not measured. These are limits of the present evaluation, not of
the design, and the open-source release includes the harness to extend it.

\paragraph{Safety and dual use.} The runtime is a planning and simulation tool for civil fleet operations. It does not
replace airworthiness assessment, airspace authorisation or onboard safety functions, and its outputs should be used
with the margins that local regulations require.
